% Generated draft: requires mathematical and source review.
\documentclass[11pt]{article}
\usepackage[margin=25mm]{geometry}
\usepackage{fontspec}
\setmainfont{Latin Modern Roman}
\usepackage{amsmath,amssymb,mathtools,unicode-math}
\setmathfont{Latin Modern Math}
\usepackage{enumitem,needspace,xurl,hyperref}
\setlength{\parindent}{0pt}
\setlength{\parskip}{5pt}
\setlength{\emergencystretch}{3em}
\newcommand{\N}{\mathbb{N}}
\newcommand{\Z}{\mathbb{Z}}
\newcommand{\Q}{\mathbb{Q}}
\newcommand{\R}{\mathbb{R}}
\newcommand{\C}{\mathbb{C}}
\newcommand{\F}{\mathbb{F}}
\newcommand{\A}{\mathbb{A}}
\newcommand{\PP}{\mathbb{P}}
\newcommand{\E}{\mathbb{E}}
\newcommand{\eps}{\varepsilon}
\newcommand{\dd}{\,\mathrm{d}}
\begin{document}
\section*{UnsolvedMath ID: 10400159}
\textbf{Source problem code:} catalogue entry
\subsection{Definitions and mathematical statement}
T. Kerler's Problem 8.13 concerns homological TQFTs. The source's collections \(\mathcal Q^*\) and \(\mathcal Q^0\), and its meanings of homological and strictly homological, are not defined in this extract. With those source meanings retained, the four requests are:
\begin{enumerate}
\item Find the irreducible components and the ring structure of \(\mathcal Q^*\) under direct sum \(\oplus\) and tensor product \(\otimes\).
\item Determine whether every strictly homological TQFT belongs to \(\mathcal Q^*\).
\item Identify the homological TQFTs coming from higher-rank gauge theories, for example \(\mathrm{PSU}(n)\) in [129], as elements of \(\mathcal Q^*\).
\item Identify the irreducible factors of the constant-order theories \(V_p^{[0]}\) in the cyclotomic-integer expansion of Reshetikhin--Turaev theory as elements of \(\mathcal Q^*\).
\end{enumerate}
The source describes (1) in terms of representation theory for the groups \(\mathrm{Sp}(2,\Z),\mathrm{Sp}(4,\Z),\ldots,\mathrm{Sp}(2g,\Z),\ldots\), together with generators and relations from standard handle attachments. It speculates that these constraints may force a positive answer to (2). For (3), it seeks understanding and possibly closed formulas for the polynomials of [129] expressing \(\mathrm{PSU}(n)\) invariants through Alexander-polynomial coefficients.

For (4), the source reports evidence suggesting \((p-3)/2\)-fold symmetric products of elements of \(\mathcal Q^0\). A proposed corollary for closed \(M\) with first Betti number \(b_1(M)\ge1\) is
\[
 V_p(M)=(\zeta_p-1)^{(p-3)/2}
 P_{(p-3)/2}\!\left(\lambda_{\rm CWL}(M)\right)
 +O\!\left((\zeta_p-1)^{(p-1)/2}\right),\tag{46}
\]
where \(\lambda_{\rm CWL}\) is the Casson--Walker--Lescop invariant, \(P_j\in\Z[x]\) has degree \(j\), and the normalization is \(V_p(S^3)=1\). The parameter \(p\), root \(\zeta_p\), coefficient ring, polynomial universality, and remainder convention must be recovered from preceding definitions; the half-integer expressions require an appropriate odd-index range. In a cyclotomic-adic interpretation the \(O\)-term means a divisibility/filtration statement, not an ordinary numerical asymptotic unless so defined. The source says [212] establishes (46) for \(p=5\) and general \(b_1\ge1\), and reports work in progress for general \(p\) when \(M\) is a torus bundle over a circle. These are source-era reports, not newly verified results.

The contextual perturbative program views homological theories as a starting point extending finite-type theory for homology \(3\)-spheres, and at fixed \(p\) views Reshetikhin--Turaev theory as a deformation of \(\mathcal Q^*\)-theories. For a surface \(\Sigma_g\), let \(\Gamma_g\) be its mapping class group and \(\mathcal I_g\) the Torelli subgroup, the kernel of the action on \(H_1(\Sigma_g;\Z)\). The action
\[
 V^{[g]}:\Gamma_g\to\mathrm{GL}_R(V(\Sigma_g))
\]
extends linearly to \(\Z[\Gamma_g]\to\operatorname{End}_R(V(\Sigma_g))\). Write \(J_g\) for the augmentation ideal of \(\Z[\mathcal I_g]\), as embedded in \(\Z[\Gamma_g]\), subject to checking whether the source instead means its generated ambient ideal. The excerpt defines length as the \emph{maximal} \(L\in\N\) satisfying
\[
 V^{[g]}(J_g^{L+1})=0.
\]
This definition has a monotonicity defect as written: annihilation of one power implies annihilation of every higher power, so it cannot specify a finite maximal \(L\). It is retained as a source issue rather than silently changed to a minimal exponent. Genus-uniformity is also unstated. The source identifies length-zero theories with homological ones and suggests length-one theories are extension objects \(\operatorname{Ext}(V,W)\), \(V,W\in\mathcal Q^*\). It further reports, in characteristic different from \(2\), factorization of the corresponding mapping-class representations through a Johnson--Morita homomorphism displayed as
\[
 \Gamma_g\longrightarrow
 \bigwedge^3 H_1(\Sigma_g)\rtimes\mathrm{Sp}(2g,\Z),
\]
with explicit extensions referenced in [211]. Coefficients, closed-versus-bordered surface conventions and the precise target of this homomorphism require source review.
\subsection{Short English statement}
Classify the source's homological TQFTs, relate gauge-theoretic and constant cyclotomic terms to them, and develop the indicated perturbative extension theory. Preserve the speculative symmetric-power and Casson--Walker--Lescop formula separately from source-reported cases. Several defining objects are missing, and the printed maximal-length definition is internally problematic.
\subsection{Sources}
\begin{itemize}
\item Local UnsolvedMath registry, permanent ID \(10400159\), exact source code \texttt{catalogue entry}, Problem 8.13, ``Homological TQFTs,'' attributed to T. Kerler. The title is truncated after ``(T.''
\item Ohtsuki, \emph{Problems on invariants of knots and 3-manifolds} (2002), Problem 8.13 and following context, PDF page 127: \url{https://msp.org/gtm/2002/04/gtm-2002-04-024s.pdf}. Bold headings distinguish 231 numbered items from surrounding references.
\item Supplied internal references: [129] for higher-rank gauge theories and Alexander-coefficient formulas; [212] for the \(p=5\) case of (46); [211] for explicit extension representations. Bibliography entries are not included. General-\(p\) torus-bundle work is described as work in progress without a separate citation.
\item Extraction pdf-font-structure, accessed 2026-07-29. Source evidence NEEDS\_REVIEW presents an open problem and requires a current primary-literature audit. Supplied research classification OPEN-TRIAGE calls Kerler's program not obviously resolved.
\item Dataset attribution: Ulam AI / UnsolvedMath Contributors, CC BY 4.0. Source rights NEEDS\_REVIEW: publicly accessible author/publisher copy, redistribution terms requiring review.
\end{itemize}
\end{document}
